\documentclass[10pt,a4paper]{article}
\usepackage{amsmath,amssymb,amsthm,graphicx}
\usepackage{amsmath,amssymb,amsthm}
\newtheorem{theorem}{Theorem}
\newtheorem{definition}[theorem]{Definition}
\newtheorem{lemma}[theorem]{Lemma}
\newtheorem{claim}[theorem]{Claim}
\newtheorem{corollary}[theorem]{Corollary}
\DeclareMathOperator{\Supp}{Supp}
\DeclareMathOperator{\CBA}{CBA}
\DeclareMathOperator{\Raz}{Raz}
\DeclareMathOperator{\Disp}{Disp}
\newcommand{\twoExt}{\mathsf{2Ext}}
\newcommand{\E}{\mathbb E}
\newcommand{\HH}{\widetilde H_\infty}
\title{Near-optimal nonlinear erasure list-decodable codes via one-bit strong dispersers and an unbalanced two-source extractor}
\author{Avraham Ben-Aroya, Dean Doron, Amnon Ta-Shma}
\date{doi:10.4230/LIPIcs.CCC.2020.1; CC BY 3.0 Unported}
\begin{document}
\maketitle
\noindent\textit{Editorial note.} Editable transcription from the official CC BY3.0 Unported publication, not relicensed4.0. Full core author proof retained: Sections2–4 plus5.1/5.3; historical comparison tables, nonstrong sketch5.2, Ramsey applications5.4 and open questions omitted. Main statement is originalTheorem2; detailedTheorems27/40/46 and all source-local supporting proofs remain in one item. Original formulas/inequalities, including Definition24 notation and Lemma44 strictness factors, are transcribed without mathematical repair or added assumptions. Only line wrapping and presentation macros adapted. Original numbered bracket references retain complete45-entry ordinary bibliography. All production/validation on Yupeng; root owns compilation.\par\medskip
\noindent\textit{Allowed standard prerequisites (editorial).}\begin{itemize}
\item Statistical distance and min-entropy, Raz strong two-source extractor, curve mergers, correlation breakers with advice, resilient functions. OriginalSection2 retains exact source-used definitions and cited statements; published standard components are cited rather than recursively reproved
\end{itemize}\medskip
\section*{Main guarantee and conventions (original Sections 1--2)}
A code $\mathcal C\subseteq\{0,1\}^{\bar n}$ is $(s,L)$ erasure list-decodable if for every word $w$, after erasing any $s$ symbols of $w$, the remaining $\bar n-s$ symbols have at most $L$ possible completions into a codeword of $\mathcal C$.
\setcounter{theorem}{1}

\begin{theorem}[see also Theorem 46]
For every constant $0<\gamma<1$ and $\varepsilon=n^{-\Omega_\gamma(1)}$ there exists an explicit $((1-\varepsilon)\bar n,L=\log^{O_\gamma(1)}\frac1\varepsilon)$ erasure list-decodable code $\mathcal C:\{0,1\}^n\to\{0,1\}^{\bar n}$ of rate $\varepsilon^{1+\gamma}$.
\end{theorem}

We stress that the codes we present are explicit in the sense that they have explicit encoding, but we do not know whether the codes we construct admit efficient erasure list-decoding algorithms. We also mention that the list-size $\operatorname{poly}(\log\frac1\varepsilon)$ achieved by our code is exponentially smaller than the best possible list-size by any linear code.
Throughout the paper we use the convention that lowercase variables are the logarithm (in base 2) of their corresponding uppercase variables, e.g., $n=\log N$, $d=\log D$. We denote by $[A]$ the set $\{1,\ldots,A\}$. The density of a set $B\subseteq A$ is $\rho(B)=|B|/|A|$. We say a function $f:A\to B$ is explicit if there exists a deterministic polynomial algorithm that runs in time $\operatorname{poly}(\log|A|)$ and computes $f$.
\section*{Preliminaries (original Section 2)}
\subsection*{Random variables and min-entropy}
The statistical distance between two distributions $X$ and $Y$ on the same domain $\Omega$ is defined as $|X-Y|=\max_{A\subseteq\Omega}(\Pr[X\in A]-\Pr[Y\in A])$. If $|X-Y|\le\varepsilon$ we say $X$ is $\varepsilon$-close to $Y$ and denote it by $X\approx_\varepsilon Y$. We denote by $U_n$ the random variable distributed uniformly over $\{0,1\}^n$. We say a random variable is flat if it is uniform over its support.
For a function $f:\Omega_1\to\Omega_2$ and a random variable $X$ distributed over $\Omega_1$, $f(X)$ is the random variable distributed over $\Omega_2$ obtained by choosing $x$ according to $X$ and computing $f(x)$. For a set $A\subseteq\Omega_1$, $f(A)=\{f(x)\mid x\in A\}$. For every $f:\Omega_1\to\Omega_2$ and two random variables $X$ and $Y$ distributed over $\Omega_1$, it holds that $|f(X)-f(Y)|\le|X-Y|$.
The min-entropy of a random variable $X$ is defined by
\[H_\infty(X)=\min_{x\in\Supp(X)}\log\frac1{\Pr[X=x]}.\]
A random variable $X$ is an $(n,k)$ source if $X$ is distributed over $\{0,1\}^n$ and has min-entropy at least $k$. When $n$ is clear from the context we sometimes omit it and simply say that $X$ is a $k$-source. Every $k$-source $X$ can be expressed as a convex combination of flat distributions each with min-entropy at least $k$.
\setcounter{theorem}{5}
\begin{definition}[average conditional min-entropy]
Let $X,Y$ be two random variables. The average conditional min-entropy of $X$ given $Y$ is
\[\HH(X\mid Y)=-\log\left(\E_{y\sim Y}\left[2^{-H_\infty(X\mid Y=y)}\right]\right).\]
\end{definition}
We will use the following simple claim about average conditional min-entropy:
\begin{claim}For any random variables $X,Y$, $\HH(X\mid Y)\ge H_\infty(X)-\log|\Supp(Y)|$.\end{claim}
\subsection*{Condensers}
\begin{definition}[condenser]
A function $C:\{0,1\}^n\times\{0,1\}^d\to\{0,1\}^m$ is an $(n,k)\to_{\varepsilon_{\rm cond}}(m,k')$ condenser, if for every $(n,k)$ source $X$, $C(X,U_d)$ is $\varepsilon_{\rm cond}$-close to an $(m,k')$ source. If $k=\delta n$ and $k'=\delta'm$ we say $C$ is a $\delta\to_{\varepsilon_{\rm cond}}\delta'$ condenser.
\end{definition}
\begin{lemma}
Suppose $C:\{0,1\}^n\times\{0,1\}^d\to\{0,1\}^m$ is an $(n,k)\to_{\varepsilon_{\rm cond}}(m,k'+d)$ condenser. Let $X$ be an $(n,k)$ source. Let $\varepsilon_i$ be the minimal distance of $C(X,i)$ to an $(m,k')$ source. Then, $\E_{i\in\{0,1\}^d}[\varepsilon_i]\le\varepsilon_{\rm cond}$.
\end{lemma}
\begin{proof}
Fix an $(n,k)$ source $X$. For $i\in\{0,1\}^d$, let $H_i\subseteq\{0,1\}^m$ be the set of ``heavy'' elements of $C(X,i)$,
\[H_i=\{w\in\{0,1\}^m:\Pr_{x\in X}[C(x,i)=w]\ge2^{-k'}\}.\]
The distance of $C(X,i)$ from a $k'$-source is
\[\varepsilon_i=\Pr_{x\in X}[C(x,i)\in H_i]-2^{-k'}|H_i|,\]
by redistributing the mass of the heavy elements. Let $H=\bigcup_{i\in\{0,1\}^d}H_i$. Then,
\begin{itemize}
\item For every $w\in H$, $\Pr_{x\in X,i\in\{0,1\}^d}[C(x,i)=w]\ge2^{-d}2^{-k'}=2^{-(k'+d)}$, and,
\item it holds that
\begin{align*}
\varepsilon_{\rm cond}&=\Pr_{x\in X,i\in\{0,1\}^d}[C(x,i)\in H]-|H|2^{-(k'+d)}\\
&=\sum_{i\in\{0,1\}^d}2^{-d}\Pr_x[C(x,i)\in H]-|H|2^{-(k'+d)}\\
&\ge\sum_{i\in\{0,1\}^d}2^{-d}\Pr_x[C(x,i)\in H_i]-2^{-(k'+d)}\sum_{i\in\{0,1\}^d}|H_i|\\
&=\sum_{i\in\{0,1\}^d}2^{-d}(\Pr_x[C(x,i)\in H_i]-2^{-k'}|H_i|)\\
&=\sum_{i\in\{0,1\}^d}2^{-d}\varepsilon_i=\E_{i\in\{0,1\}^d}[\varepsilon_i].
\end{align*}
\end{itemize}
\end{proof}
\subsection*{Two-source extractors}
\begin{definition}[two-source extractor]
A function $\twoExt:\{0,1\}^{n_1}\times\{0,1\}^{n_2}\to\{0,1\}^m$ is an $((n_1,k_1),(n_2,k_2),\varepsilon)$ two-source extractor if for every two independent sources $X_1$ and $X_2$ where $X_1$ is an $(n_1,k_1)$ source and $X_2$ is an $(n_2,k_2)$ source, it holds that $\twoExt(X_1,X_2)\approx_\varepsilon U_m$. We say that $\twoExt$ is strong if
\[(\twoExt(X_1,X_2),X_1)\approx_\varepsilon(U_m,X_1)\quad\text{and}\quad(\twoExt(X_1,X_2),X_2)\approx_\varepsilon(U_m,X_2).\]
\end{definition}
In our construction, we will use the following two-source extractor:
\begin{theorem}[{[39]}]
For every constant $\delta_{\rm Raz}>\frac12$ there exist constants $c_1=c_1(\delta_{\rm Raz}),c_2=c_2(\delta_{\rm Raz})>1$ such that for every $n_1,k_1,n_2,k_2$ satisfying
\[k_1\ge c_1\log n_2,\qquad k_2\ge c_2\log n_1,\]
there exists an explicit strong $((n_1,k_1),(n_2,k_2=\delta_{\rm Raz}n_2),\varepsilon_{\rm Raz})$ two-source extractor
\[\Raz:\{0,1\}^{n_1}\times\{0,1\}^{n_2}\to\{0,1\}^m\]
with $m=\Omega(\min\{k_1,k_2\})$ and $\varepsilon_{\rm Raz}=2^{-\Omega(m)}$, where the constants hiding in the asymptotic notation may depend on $\delta_{\rm Raz}$.
\end{theorem}
\begin{claim}
Suppose $\twoExt:\{0,1\}^{n_1}\times\{0,1\}^{n_2}\to\{0,1\}^m$ is a strong $((n_1,k_1),(n_2,k_2),\varepsilon)$ two-source extractor. Let $X$ be an $(n,k_1)$ source. Call an element $y\in\{0,1\}^{n_2}$ bad if $|\twoExt(X,y)-U_m|>\varepsilon$, and let $B_Y$ denote the set of all bad elements. Then, $|B_Y|<2^{k_2}$.
\end{claim}
\begin{proof}
Assume towards contradiction that $|B_Y|\ge2^{k_2}$ and let $Y$ be the uniform distribution over the set $B_Y$. Then, $H_\infty(Y)\ge k_2$ and so $(\twoExt(X,Y),Y)\approx_\varepsilon(U_m,Y)$ which implies that
\[\frac1{|B_Y|}\sum_{y\in B_Y}|\twoExt(X,y)-U_m|\le\varepsilon.\]
However, $|\twoExt(X,y)-U_m|>\varepsilon$ for every $y\in B_Y$, in contradiction.
\end{proof}
\subsection*{Mergers}
A merger takes as input a list of possibly correlated random variables along with a short uniform seed and outputs one random variable which is close to having high min-entropy, provided at least one of the input variables has high min-entropy. Formally:
\begin{definition}[somewhere-random source]
A source $X=X_1\circ\ldots\circ X_A$ is an $(n,k,(\alpha,\beta))$ somewhere-random (s.r.) source if there is a random variable $I\in\{0,\ldots,A\}$ such that for every $i\in[A]$, $(X_i\mid I=i)$ is $\alpha$-close to an $(n,k)$ source and $\Pr[I=0]\le\beta$. The variable $I$ is called the indicator of source. If $\alpha=\beta=0$ we say $X$ is a $(n,k)$ s.r. source.
\end{definition}
\begin{definition}[merger]
A function $B:(\{0,1\}^n)^D\times\{0,1\}^t\to\{0,1\}^m$ is a $(k,k',\varepsilon)$ merger, if for every $(n,k)$ s.r. source $X=X_1\circ\ldots\circ X_A$, the output $M(X,U_t)$ is $\varepsilon$-close to a $k'$-source.
\end{definition}
There are explicit constructions of good mergers. Dvir and Wigderson [17] constructed the curve merger and proved that it works with $t=O(\log\frac n\varepsilon)$. This was further improved in [16] who proved that $t=O(\log\frac D\varepsilon)$ suffices. Notice that now $t$ only depends on the number of sources $D$ and the requested error $\varepsilon$, but not on the source length $n$, and this remarkable property will be crucial for us. Formally,
\begin{theorem}[{[17,16]}]
There exists a constant $c_{\rm DKSS}\ge1$ such that the following holds. Fix $\beta,\delta,\varepsilon>0$. There exists an explicit function $B:(\{0,1\}^n)^D\times\{0,1\}^t\to\{0,1\}^n$ that is a $(k=\delta n,k'=(1-\beta)\delta n,\varepsilon)$ merger, with $t=c_{\rm DKSS}\cdot\frac1\beta\log\frac D\varepsilon$.
\end{theorem}
\subsection*{Correlation breakers with advice}
A correlation-breaker with advice is a function $\CBA:\{0,1\}^n\times\{0,1\}^\ell\times[A]\to\{0,1\}^m$ where we think of the first input as a weak source, the second as an independent short seed and the last as an advice string. Roughly speaking, applying CBA on $t$ possibly correlated seeds with $t$ distinct advice strings results in independent random variables. For example, $\CBA(X,Y,\alpha)$ is (nearly) independent of $\CBA(X,Y,\alpha')$ for any $\alpha\ne\alpha'$. Formally,
\begin{definition}
A function $\CBA:\{0,1\}^n\times\{0,1\}^\ell\times[A]\to\{0,1\}^m$ is a $(t,k,\varepsilon_{\rm CBA})$ correlation-breaker with advice if the following holds. If $Y$ is a distribution over $\{0,1\}^n$, $Z=(Z_1,\ldots,Z_t)$ is a distribution on $(\{0,1\}^\ell)^t$, $H$ is a random variable and $\delta>0$ satisfy:
\begin{itemize}
\item $Y$ and $Z$ are independent, conditioned on $H$,
\item $\HH(Y\mid H)\ge k+\log(1/\varepsilon_{\rm CBA})$,
\item $(Z_1,H)\approx_\delta(U_\ell,H)$, and,
\item $\alpha_1,\ldots,\alpha_t\in[A]$ are distinct strings.
\end{itemize}
Then,
\begin{multline*}
(\CBA(Y,Z_1,\alpha_1),(\CBA(Y,Z_i,\alpha_i))_{i=2}^t,H)\\
\approx_{\delta+2\varepsilon_{\rm CBA}}(U_m,(\CBA(Y,Z_i,\alpha_i))_{i=2}^t,H).
\end{multline*}
\end{definition}
We use the following result:
\begin{theorem}[{[13, Theorem 4.12]}]
There exists a constant $c_{\rm CBA}\ge1$ such that the following holds. Let $n,a$ be integers and $\varepsilon_{\rm CBA}>0$. Then, there exists an explicit $(t,k_{\rm CBA},\varepsilon_{\rm CBA})$ correlation-breaker with advice
\[\CBA:\{0,1\}^n\times\{0,1\}^\ell\times[A]\to\{0,1\}\]
with $\ell=c_{\rm CBA}\cdot at\cdot\log\frac n{\varepsilon_{\rm CBA}}$ and $k_{\rm CBA}\ge\ell$.
\end{theorem}
In our setting, the number of rows $A$ is a constant independent of $n$. For this reason we work with a ``basic'' correlation-breaker, where there is no attempt to optimize the dependence of $\ell$ on $a$. This gives a seed-length which is optimal up to constant multiplicative factors.
We also need the following lemma.
\begin{lemma}
Let $X_1,\ldots,X_t$ be random variables over $\{0,1\}^m$. Further suppose that for any $i\in[t]$,
\[(X_i,\{X_j\}_{j\ne i})\approx_\varepsilon(U_m,\{X_j\}_{j\ne i}).\]
Then, $(X_1,\ldots,X_t)\approx_{t\varepsilon}U_{tm}$.
\end{lemma}
\subsection*{Limited independence and non-oblivious bit-fixing sources}
\begin{definition}
A distribution $X$ over $\{0,1\}^A$ is called $(t,\gamma)$-wise independent if the restriction of $X$ to every $t$ coordinates is $\gamma$-close to $U_t$. A source $X$ over $\{0,1\}^A$ is called a $(q,t,\gamma)$ non-oblivious bit-fixing source if there exists a subset $Q\subseteq A$ of size at most $q$ such that the joint distribution of the bits in $A\setminus Q$ is $(t,\gamma)$-wise independent. The bits in $Q$ are allowed to arbitrarily depend on the bits in $A\setminus Q$. If $\gamma=0$ we often say that $X$ is a $(q,t)$ non-oblivious bit-fixing source.
\end{definition}
\begin{lemma}[{[3]}]A $(t,\gamma)$-wise distribution over $A$ bits is $(A^t\gamma)$-close to some $t$-wise independent distribution.\end{lemma}
\begin{definition}
Let $f:\{0,1\}^A\to\{0,1\}$, $\mathcal D$ a distribution over $\{0,1\}^A$ and $Q\subseteq A$. Let $I_{Q,\mathcal D}(f)$ denote the probability that $f$ is undetermined when the variables outside $Q$ are sampled from $\mathcal D$. We define $I_{q,t,\gamma}(f)$ to be the maximum of $I_{Q,\mathcal D}(f)$ over all $Q\subseteq A$ of size $q$ and all $\mathcal D$ that is a $(t,\gamma)$-wise independent distribution. We say that $f$ is $(t,\gamma)$-independent $(q,\varepsilon)$-resilient if $I_{q,t,\gamma}(f)\le\varepsilon$.
\end{definition}
\begin{theorem}[{[9,32]}]
For every $0<\gamma<1$ there exists a constant $c_\gamma\ge1$ such that for all $A>0$ there exists an explicit function $f:\{0,1\}^A\to\{0,1\}$ with the following property: For every $t\ge c_\gamma\log^4 A$,
\begin{itemize}
\item $f$ is almost balanced: For any $t$-wise independent distribution $\mathcal D$ on $\{0,1\}^A$, $\Pr_{x\sim\mathcal D}[f(x)=1]=1/2\pm A^{-1/c_\gamma}$, and,
\item $f$ is resilient: $I_{q,t,\gamma}(f)\le c_\gamma\cdot q/A^{1-\gamma}$.
\end{itemize}
\end{theorem}

\section*{Constant Degree Condensers (original Section 3)}
In this section we prove:
\begin{theorem}
For every constant $0<\delta_1<\delta_2=0.7$, every $s\ge s_0(\delta_1)$ and every integer $n_1$ and $\varepsilon_{\rm cond}\ge2^{-\Omega(n_1)}$ there exists an explicit $\delta_1\to_{\varepsilon_{\rm cond}}\delta_2$ condenser $C:\{0,1\}^{n_1}\times\{0,1\}^d\to\{0,1\}^{n_2}$ with $n_2=(\frac23)^sn_1$ and $d=4c_{\rm DKSS}(s+\log\frac1{\varepsilon_{\rm cond}})$, where $c_{\rm DKSS}$ is the constant from Theorem 15. Note that $d$ is independent of $n_1$.
\end{theorem}
Note that, in particular, for every $\delta_1>0$ there exists an explicit $\delta_1\to_{\varepsilon_{\rm cond}}\delta_2=0.7$ condenser $C:\{0,1\}^{n_1}\times\{0,1\}^d\to\{0,1\}^{n_2}$ with $n_2=\Omega(n_1)$ and $d=O(\log\frac1{\varepsilon_{\rm cond}})$. However, we will need the more precise version that appears in Theorem 23.
The proof goes through somewhere-random condensers, so let us first discuss the similarities and differences between condensers and somewhere-random condensers. We begin with the necessary definitions:
\begin{definition}[s.r. condenser]
A function $C:\{0,1\}^n\to(\{0,1\}^m)^A$ is an $(n,k)\to_\varepsilon(m,k')$ s.r. condenser if for every $(n,k)$ source $X$ it holds that $C(X)=C(X,1)\circ\ldots\circ C(X,A)$ is $\varepsilon$-close to a $k'$ s.r. source. If $k=\delta n$ and $k'=\delta m$ we say $C$ is $\delta\to_\varepsilon\delta'$ s.r. condenser.
\end{definition}
We may take a condenser $C:\{0,1\}^n\times\{0,1\}^d\to\{0,1\}^m$ and expand it to a table with the outputs of all possible seeds, i.e., define $S:\{0,1\}^n\to(\{0,1\}^m)^D$, with $D=2^d$, where $S(x)_i=C(x,i)$. The condenser property guarantees that for every $k$-source $X$, most rows in the table are close to having $k'$ min-entropy. In contrast, a s.r. condenser is a weaker object, because it only guarantees that one row has $k'$ entropy (or more precisely that we are in a convex combination of such cases).
The major question we consider now is the dependence of the degree ($2^d$ for condensers and $A$ for s.r. condensers) on $n,m,k,k'$ and $\varepsilon$. We focus on the case where $m=\Omega(n)$, $k=\delta n$, $k'=\delta'm$ and $\delta<\delta'$ are constants. A priori, we could have expected the degree to depend on $n$ and $\varepsilon$, as is indeed the case when $m$ might be arbitrarily small. However, remarkably, things are drastically different when $m=\Omega(n)$. In this case both condensers and s.r. condensers may be of degree that is independent of $n$ and this will be crucial for us. If we consider the dependence on the error, then s.r. condensers may have exponentially-small error and constant $D$, whereas the degree of a condenser is at least $d\ge\log(\frac1\varepsilon)$. Remarkably, all of that can be explicitly achieved, as we now explain.
The basic building block we use is the following beautiful result of Zuckerman, which is based on additive combinatorics:
\begin{theorem}[{[45, Theorem 8.3]}]
For every constant $0<c<1$ there exists a constant $\alpha=\alpha(c)$ such that for every constant $\delta\le c$ and integer $n$ there exists an explicit function $C:\{0,1\}^n\to(\{0,1\}^{\frac23n})^2$ that is a $\delta\to_\varepsilon(1+\alpha)\delta$ s.r. condenser with $\varepsilon=2^{-\Omega(\alpha\delta n)}$.
\end{theorem}
Somewhere-random condensers can be easily composed. Specifically, Barak et al. [5] showed that if $C_1:\{0,1\}^{n_1}\to(\{0,1\}^{n_2})^{\ell_1}$ is a $\delta_1\to_{\varepsilon_1}\delta_2$ s.r. condenser and $C_2:\{0,1\}^{n_2}\to(\{0,1\}^{n_3})^{\ell_2}$ a $\delta_2\to_{\varepsilon_2}\delta_3$ s.r. condenser then $C_2\circ C_1:\{0,1\}^{n_1}\to(\{0,1\}^{n_3})^{\ell_1\cdot\ell_2}$ defined by $C_2\circ C_1(x)_{(i_1,i_2)}=C_2(C_1(x)_{i_1})_{i_2}$ is a $\delta_1\to_{\varepsilon_1+\varepsilon_2}\delta_3$ s.r. condenser.
Composing the s.r condenser of Theorem 25 with itself $s$ times we get an explicit function $C:\{0,1\}^n\to(\{0,1\}^m)^D$ with $D=2^s$ and $m=(\frac23)^sn$ that is a $\delta\to_\varepsilon\delta'$ s.r. condenser with $\varepsilon=\sum_{i=1}^s2^{-\Omega((1+\alpha)^i\delta(\frac23)^in)}=2^{-\Omega(m)}$ and $\delta'\ge(1+\alpha(\delta'))^s\delta$. Therefore:
\begin{lemma}
For every constants $0<\delta_1<\delta_2<1$ there exists a constant $s=s(\delta_1,\delta_2)$ and an explicit function $C:\{0,1\}^{n_1}\to(\{0,1\}^{n_2})^D$ that is a $\delta_1\to_\varepsilon\delta_2$ s.r. condenser with $D=2^s$, $n_2=(\frac23)^sn_1$ and $\varepsilon=2^{-\Omega(n_2)}$. Note that $D$ is independent of $n$ and $\varepsilon$.
\end{lemma}
Right now, if $X$ is a $k$-source, the table $C(X)$ has $D$ rows, and, roughly speaking, the guarantee is that one of these rows has density $\delta'$. We want to change this to get a condenser, i.e., we are willing to invest a short seed (that is independent of $n$) and we want to get one output which is close to uniform. (Alternatively, we can write the condenser as a table with one row per seed, the number of rows is independent of $n$ and most rows are close to uniform.) This is exactly what a merger does and applying the merger of Theorem 15 with $\beta=\frac14$ on the s.r. condenser of Lemma 26 (with $\delta_2$ close to 1) gives Theorem 23.
\section*{The Unbalanced Two-Source Extractor Construction (original Section 4)}
The main result of this section is the following two-source extractor.
\begin{theorem}
For every integer $n$ and two constants $\delta_0,\varepsilon_0>0$ there exists a constant $c$ such that for $d\ge c\log n$ and $k\ge c\log d$ there exists an explicit $((n,k),(d,\delta_0d),\varepsilon_0)$ two-source extractor $\twoExt:\{0,1\}^n\times\{0,1\}^d\to\{0,1\}$.
\end{theorem}
The extractor in the above theorem has constant error, and works when:
\begin{enumerate}
\item Each source's entropy is in the order of the logarithm of the length of the other source.
\item The shorter source, of length $d$, has an arbitrarily small constant density $\delta_0$.
\end{enumerate}
We think of $n$ and $d=d(n)$ as growing parameters while $\varepsilon_0$ and $\delta_0$ are constants. We use asymptotic notations (such as $\Omega(\cdot)$) to hide constants that are independent of $n$ and $d$ (but may depend on $\varepsilon_0$ and $\delta_0$).
\subsection*{4.1 The Construction}
Recall that $\varepsilon_0$ is the target error of the extractor $\twoExt$. The input to $\twoExt$ is a pair $(x,y)$ where $x$ is drawn from an $(n,k)$ source $X$, and $y$ is drawn from an independent $(d,\delta_0d)$ source $Y$. Our problem is that the $y$ comes from a $\delta_0d$-source for some $\delta_0<\frac12$. To overcome this, we do the following:
\begin{itemize}
\item We apply the condenser of Theorem 23 on $y$ to get a table $y'$ that is 1-wise 0.7-dense. Notice that the output of this step is a table rather than a single output.
\item We apply Raz's extractor (Theorem 11) on the table and the input $x$ from the other source to convert the table $y'$ to another table $y''$ that is 1-wise uniform.
\item We apply the $t$ correlation-breaker with advice of Theorem 17 on $y$, using the table $y''$ as the seed, to get a table $y'''$ that is $t$-wise uniform.
\item Finally, we apply the resilient function $f$ of Theorem 22 on the table $y'''$ to collapse the many rows of the table to a single, close to uniform, output.
\end{itemize}
Formally, these steps work as follows:
\paragraph{Condense the short source.} We are given $\delta_0<\frac12$. Set $\delta'=0.69$ and $\delta_2=0.7$.
By Theorem 23 there exists a constant $s_0=s_0(\delta_0)$ such that for every $s\ge s_0$ there exists an explicit
\[C:\{0,1\}^d\times\{0,1\}^a\to\{0,1\}^{d'}\]
that is a $\delta_0\to_{\varepsilon_{\rm cond}}\delta_2=0.7$ condenser with $a=4c_{\rm DKSS}(s+\log\frac1{\varepsilon_{\rm cond}})$ and $d'=(\frac23)^sd$.
We set
\[\gamma=\frac1{2^5c_{\rm DKSS}},\]
and this also fixes $c_\gamma$ as in Theorem 22. Notice that $\gamma$ and $c_\gamma$ are fixed constants independent of all other parameters in our system.
Now, choose $\varepsilon_{\rm cond}$ so that
\begin{equation}\left(\frac1{\varepsilon_{\rm cond}}\right)^{\log(3/2)}\ge\frac4{\delta_0}2^{12}c_\gamma c_{\rm DKSS}^4\log^4\frac1{\varepsilon_{\rm cond}},\tag{1}\label{eq:cond}\end{equation}
and also so that $\varepsilon_{\rm cond}\le\xi(\varepsilon_0,\delta_0)$, where
\begin{equation}\xi(\varepsilon_0,\delta_0)=\min\left\{2^{-s_0},\left(\frac{\varepsilon_0}8\right)^2,\left(\frac{\varepsilon_0}5\right)^{c_\gamma},\left(\frac{\varepsilon_0}{5c_\gamma}\right)^{1/\gamma}\right\}.\tag{2}\label{eq:xi}\end{equation}
Given $\varepsilon_{\rm cond}$, we set $s=\log\frac1{\varepsilon_{\rm cond}}\ge s_0$, giving $a=8c_{\rm DKSS}\log\frac1{\varepsilon_{\rm cond}}$. Note that the degree of the condenser, $A=2^a$, satisfies
\[\sqrt{\varepsilon_{\rm cond}}A=2^{-\frac12\log\frac1{\varepsilon_{\rm cond}}+a}=2^{-\frac a{2^4c_{\rm DKSS}}+a}=A^{1-2\gamma}.\]
Observe that $s\ge s_0$ and that $d'=\Omega(d)$. Also, notice that $(\delta_2-\delta')d'=d'/100\ge a=\log A$ for large enough $d$. Thus, $C$ is a $(d,\delta d)\to_{\varepsilon_{\rm cond}}(d',\log(A)+\delta'd')$ condenser.
Define an $A\times d'$ table $Y'$ where
\[Y'_i=C(Y,i)\in\{0,1\}^{d'}\quad\text{for }i=1,\ldots,A.\]
\paragraph{1-wise uniformity.} Let $c_1,c_2$ be the constants from Theorem 11 for $\delta_{\rm Raz}=0.6$.
Notice that $\delta_{\rm Raz}d'=\Omega(d')=\Omega(d)$. Therefore, for a constant $c$ large enough, $d\ge c\log n$ is large enough so that $\delta_{\rm Raz}d'\ge c_2\log n$. We can, in particular, choose $c$ such that in addition $c\ge c_1$. Recalling that $k\ge c\log d$, we have $k\ge c_1\log d'$. By Theorem 11, there exists an explicit function
\[\Raz:\{0,1\}^n\times\{0,1\}^{d'}\to\{0,1\}^{d''}\]
that is a strong $((n,k),(d',\delta_{\rm Raz}d'),\varepsilon_{\rm Raz}=2^{-\Omega(d'')})$ two-source extractor with $d''=\Omega(\min\{k,\delta_{\rm Raz}d'\})=\Omega(k)$.\footnote{Although $k\ge c\log d$ we can always assume w.l.o.g. that $k=c\log d$ and so $k\le\delta_{\rm Raz}d'=\Omega(d)$.} Define an $A\times d''$ table $Y''$ where
\[Y''_i=\Raz(X,Y'_i)\quad\text{for }i=1,\ldots,A.\]
\paragraph{$t$-wise uniformity.} Let $k_{\rm CBA}=\frac{\delta_0d}8$ and $\varepsilon_{\rm CBA}=\frac1d$. Set
\[t=\frac{\delta_0}4\left(\frac32\right)^s.\]
Notice that for a large enough constant $c$ we have $d''=\Omega(k)=\Omega(c\log d)\ge c_{\rm CBA}at\log\frac d{\varepsilon_{\rm CBA}}$, where the latter is the seed-length required by the correlation-breaker from Theorem 17. Also, $k_{\rm CBA}=\frac{\delta_0d}8\ge d''$ for large enough $d$, as $d''=\Omega(k)=\Omega(\log d)$. Hence, by Theorem 17 there exists an explicit function
\[\CBA:\{0,1\}^d\times\{0,1\}^{d''}\to\{0,1\}\]
that is a $(t,k_{\rm CBA},\varepsilon_{\rm CBA})$ correlation-breaker with advice.
Define an $A\times1$ table $Y'''$ where
\[Y'''_i=\CBA(Y,Y''_i,i)\quad\text{for }i=1,\ldots,A.\]
Keep in mind that the entropy in $Y$ suffices for CBA since $H_\infty(Y)=8k_{\rm CBA}$.
\paragraph{Collapse.} Take $f:\{0,1\}^A\to\{0,1\}$ to be the $(q=A^{1-2\gamma},t,\varepsilon_f=c_\gamma A^{-\gamma})$ resilient function of Theorem 22 and output $f(y'''_1,\ldots,y'''_A)$.
\subsection*{4.2 Two Subtleties}
As mentioned in the introduction, there are several delicate issues in the analysis:
\begin{enumerate}
\item Circular dependence: $Y''$ depends on both $X$ and $Y$, and is used as a seed in the application of the correlation-breaker with advice on $Y$.
\item We need $Y'''$ to be close to a perfect $t$-wise independent table, while the correlation-breaker with advice only guarantees that every $t$ good rows are close to uniform. To bridge the gap we need the error to be at least polynomially-small in the number of rows, but some of the steps incur a large constant error.
\end{enumerate}
To overcome the first issue we show a conditioning under which $Y''$ is still good, $Y$ is independent of $Y''$ and even after the conditioning the two sources have enough min-entropy.
To overcome the second issue we distinguish between large errors that depend on the number of rows $A$, and small errors that depend on the row length (see Section 1.3 in the introduction of the original publication). In particular, the errors are of three types:
\begin{itemize}
\item The probability $p_1$ that a value we condition upon is bad. This error is incurred by the condenser and is high (think of it as being a constant).
\item We show that when we condition on a good value, every $t$ good rows in $Y'''$ are $p_2$-close to uniform. We then claim that $Y'''$ as a table is $A^tp_2$-close to a table where the good rows are truly $t$-wise independent (where $A$ is the number of rows in the table $Y'''$). The error $p_2$ is incurred by Raz's extractor and by the correlation-breaker, and can be made very small if we deal with a source $X$ having enough min-entropy. We make $p_2$ small enough so that $A^tp_2$ is also small.
\item A third error $p_3$ is incurred by the resilient function $f$. This error is large, say, a constant, and we are fine with that.
\end{itemize}
Note that we cannot just accumulate all errors as $A^tp_1$ is way larger than 1. Instead, we argue that with a constant probability $1-p_1$, we get extremely close to perfect behavior, and then we get such a small error $p_2$ so that $A^tp_2$ is also small.

\subsection*{4.3 The Analysis}
\paragraph{Proof of Theorem 27.} Fix an $(n,k)$ source $X$ and an independent $(d,\delta d)$ source $Y$. We decompose the proof into three parts:
\begin{itemize}
\item In the first part we prove that very often (except for a small constant probability) the table $Y''$ contains many rows that are marginally close to uniform.
\item Next, we prove that every set of $t$ rows $\{i,j_1,\ldots,j_{t-1}\}$ in $Y'''$ are product in the sense that if $i$ is a good row (intuitively meaning that $Y''_i$ is close to uniform) and $j_1,\ldots,j_{t-1}$ are $t-1$ other rows, then in $Y'''$, $Y'''_i$ is close to uniform and independent of $Y'''_{j_1},\ldots,Y'''_{j_{t-1}}$. This part involves applying a correlation-breaker with advice on $Y$ and $Y''$. In order to ensure that $Y$ and $Y''$ are independent, we condition on the values of $Y'$ in the $t$ rows $\{i,j_1,\ldots,j_{t-1}\}$.
\item Together, except for a small constant probability, there are many good rows, and every $t$ rows of $Y'''$ are product, hence the table $Y'''$ is close to a $(q,t)$ non-oblivious bit-fixing source, where every good row is a good bit in the bit-fixing source. Hence, $f(Y''')$ is close to uniform.
\end{itemize}
\paragraph{Part 1: Often, many rows in $Y'$ are good.}
Let $\varepsilon_i$ be the minimal distance of $C(Y,i)$ from a $\delta'd'$-source. According to Lemma 9,
\[\E_{i\in[A]}[\varepsilon_i]\le\varepsilon_{\rm cond}.\]
\begin{definition}
We say $z\in\{0,1\}^{d'}$ is good if $\Raz(X,z)$ is $\varepsilon_{\rm Raz}$-close to uniform. Let $G_Z$ be the set of all good $z$-s, and $B_Z$ the rest. We say $i\in[A]$ is good for $y\in\{0,1\}^d$ if $C(y,i)\in G_Z$ and bad otherwise. We define a random variable $B_i$, where the sample space is $Y$, and $B_i(y)=1$ if $i$ is bad for $y$ and 0 otherwise.
\end{definition}
By Claim 12, $|B_Z|\le2^{\delta_{\rm Raz}d'}$. Therefore, in expectation, the number of bad rows for $y$ is small:
\begin{claim}$\E_{y\in Y}[\sum_{i\in[A]}B_i(y)]\le2\varepsilon_{\rm cond}A$.\end{claim}
\begin{proof}
Fix an $i\in[A]$. We have that $C(Y,i)$ is $\varepsilon_i$-close to some $\delta'd'=0.69d'$-source $R$. Hence:
\[\E_y[B_i(y)]=\Pr_{y\in Y}[C(y,i)\in B_Z]\le\varepsilon_i+\Pr_{r\in R}[r\in B_Z]\le\varepsilon_i+\frac{|B_Z|}{2^{\delta'd'}}=\varepsilon_i+2^{-0.09d'}.\]
Thus, for $d$ large enough,
\begin{align*}
\E_y\left[\sum_{i\in[A]}B_i(y)\right]&=\sum_{i\in[A]}\E_y[B_i(y)]\le\sum_{i\in[A]}(\varepsilon_i+2^{-0.09d'})\\
&\le\varepsilon_{\rm cond}A+2^{-0.09d'}A\le2\varepsilon_{\rm cond}A.
\end{align*}
\end{proof}
\begin{definition}
We say $y\in\Supp(Y)$ has many bad rows if $\sum_{i\in[A]}B_i(y)\ge\sqrt{\varepsilon_{\rm cond}}A$.
\end{definition}
Denote $p_{1,1}=\varepsilon_0/4$.
\begin{claim}$\Pr_{y\in Y}[y\text{ has many bad rows}]\le p_{1,1}$.\end{claim}
\begin{proof}
By Markov,
\begin{align*}
\Pr_{y\in Y}\left[\sum_iB_i(y)\ge\sqrt{\varepsilon_{\rm cond}}A\right]&\le\frac{\E[\sum_iB_i(y)]}{\sqrt{\varepsilon_{\rm cond}}A}\\
&\le\frac{2\varepsilon_{\rm cond}A}{\sqrt{\varepsilon_{\rm cond}}A}=2\sqrt{\varepsilon_{\rm cond}}\le\frac{\varepsilon_0}4,
\end{align*}
where the last inequality follows from the fact that $\varepsilon_{\rm cond}\le(\varepsilon_0/8)^2$.
\end{proof}
\paragraph{Part 2: The good rows are $t$-wise independent.}
We introduce some notations to simplify the expressions in the proof. For $y_0\in\{0,1\}^d$ and $k\in[A]$, let $Y'''_k(y_0)$ denote $(Y'''_k\mid Y=y_0)$. Also, for a set $S\subseteq[A]$, define $Y'''_S(y_0)=\{Y'''_j(y_0)\}_{j\in S}$. Denote $p_2=\varepsilon_{\rm Raz}+2\varepsilon_{\rm CBA}$.
\begin{definition}
Let $y_0\in\{0,1\}^d$ (not necessarily in the support of $Y$). Let $i\in[A]$ and $S\subseteq[A]\setminus\{i\}$ of cardinality $t-1$. We say $y_0$ violates the product rule for $(i,S)$ if $B_i(y_0)=0$ and
\[(Y'''_i(y_0),Y'''_S(y_0))\not\approx_{p_2}U_1\times Y'''_S(y_0).\]
\end{definition}
\begin{definition}
Let $y_0\in\{0,1\}^d$ (not necessarily in the support of $Y$). Let $i\in[A]$ and $S\subseteq[A]\setminus\{i\}$ of cardinality $t-1$. We say $y_0$ violates the product rule with distinguisher $\Delta:\{0,1\}^t\to\{0,1\}$ for $(i,S)$ if $B_i(y_0)=0$ and
\[\left|\Pr[\Delta(Y'''_i(y_0),Y'''_S(y_0))=1]-\Pr[\Delta(U_1,Y'''_S(y_0))=1]\right|>p_2.\]
\end{definition}
Observe that if $y_0$ violates the product rule then there exists some $\Delta$ such that $y_0$ violates the product rule with distinguisher $\Delta$.
\begin{lemma}
For every $i$ and $S$ as above, the number of $y\in\{0,1\}^d$ that violate the product rule for $(i,S)$ is at most $2^{\delta_0d/2+2^t}$.\footnote{We could have used an alternative argument that avoids the $2^{2^t}$ factor here by a minor deterioration in the error of the CBA. However, since the $t$ we use is constant the $2^{2^t}$ factor is negligible.}
\end{lemma}
\begin{proof}
Suppose the lemma is false for some $(i,S)$. Then, by the pigeonhole principle there exists some $\Delta$ such that the number of elements $y\in\{0,1\}^d$ that violate the product rule for $(i,S)$ with distinguisher $\Delta$ is at least $2^{\delta_0d/2}$. Let $BY$ denote the set of these elements. Identify $BY$ with the uniform distribution over the set $BY$.
Let $BY'_i=C(BY,i)$, $BY''_i=\Raz(X,BY'_i)$ and $BY'''_i=\CBA(BY,BY''_i,i)$. For a subset $T\subseteq[A]$ Let $BY'_T$ denote the sub-table of $BY'$ corresponding to the rows of $T$, and similarly $BY''_T$ and $BY'''_T$. Since for every $y\in BY$, we have that
\[\Delta(BY'''_i(y),BY'''_S(y))\not\approx_{p_2}\Delta(U_1,BY'''_S(y)),\]
this holds also on average, that is
\[\Delta(BY'''_i,BY'''_S)\not\approx_{p_2}\Delta(U_1,BY'''_S).\]
Thus, it follows that
\begin{equation}BY'''_{S\cup\{i\}}\not\approx_{p_2}U_1\times BY'''_S.\tag{3}\label{eq:product}\end{equation}
On the other hand, when we condition on the values of $H=BY'_{S\cup\{i\}}$, the conditions for the correlation-breaker with advice hold:
\begin{itemize}
\item $BY$ and $BY''_{S\cup\{i\}}$ are independent given $H=BY'_{S\cup\{i\}}$, since $H$ is a function of $BY$ alone, and given that $H=BY'_{S\cup\{i\}}=h$ for some $h$, $BY''_{S\cup\{i\}}$ is a function of $X$ alone.
\item It holds that
\begin{align*}
\HH(BY\mid H)&\ge H_\infty(BY)-\log(|\Supp(H)|)\\
&=H_\infty(BY)-td'\ge\frac{\delta_0d}2-td'\ge\frac{\delta_0d}4,
\end{align*}
because
\[\frac{td'}d=t\cdot\left(\frac23\right)^s=\frac{\delta_0}4.\]
Now, since $k_{\rm CBA}=\frac{\delta_0d}8$ and $\varepsilon_{\rm CBA}=\frac1d$ we also have for $d$ large enough,
\[\HH(BY\mid H)\ge\frac{\delta_0d}4\ge k_{\rm CBA}+\log\frac1{\varepsilon_{\rm CBA}}.\]
\item $B_i(y)=0$, hence $BY'_i\in G_Z$ and $BY''_i=\Raz(X,BY'_i)$ is $\varepsilon_{\rm Raz}$-close to uniform.
\end{itemize}
Thus, by the correlation-breaker with advice property,
\begin{multline*}
(\CBA(BY,BY''_i,i),\{\CBA(BY,BY''_j,j)\}_{j\in S})\\
\approx_{\varepsilon_{\rm Raz}+2\varepsilon_{\rm CBA}}(U_1,\{\CBA(BY,BY''_j,j)\}_{j\in S}),
\end{multline*}
or, equivalently,
\[(BY'''_i,BY'''_S)\approx_{p_2}U_1\times BY'''_S,\]
in contradiction to Equation (3).
\end{proof}
\begin{definition}
Say $y\in\{0,1\}^d$ violates the product rule if it violates it for some $i\in[A]$ and $S\subseteq[A]\setminus\{i\}$ of cardinality $t-1$.
\end{definition}
As $H_\infty(Y)\ge\delta_0d$, the probability $y\in Y$ violates the product rule for a specific $(i,S)$ is at most $2^{\delta_0d/2+2^t-\delta_0d}=2^{2^t-\delta_0d/2}$. Let $p_{1,2}=\varepsilon_0/10$. Then, by the union bound, for $d$ large enough:
\begin{corollary}$\Pr_{y\in Y}[y\text{ violates the product rule}]\le2^{2^t-\delta_0d/2}\cdot A^t\le p_{1,2}$.\end{corollary}
\paragraph{Part 3: Completing the proof.}
\begin{definition}
We say $y$ is bad if it has many bad rows or if it violates the product rule. If $y$ is not bad we say it is good.
\end{definition}
Let $p_1=p_{1,1}+p_{1,2}$. Clearly, by Claim 31 and Corollary 36, $\Pr_{y\in Y}[y\text{ is bad}]\le p_1=(\frac14+\frac1{10})\varepsilon_0$.
\begin{claim}
Fix any good $y\in Y$. Then, $Y'''(y)$ is a $(q,t,tp_2)$ non-oblivious bit-fixing source, for $q=\sqrt{\varepsilon_{\rm cond}}A$.
\end{claim}
\begin{proof}
Let $Q(y)\subseteq[A]$ be the set of bad rows for $y$. As $y$ does not have many bad rows, $|Q(y)|=\sum_{i\in[A]}B_i(y)\le\sqrt{\varepsilon_{\rm cond}}A=q$.
Now, fix any set $S\subseteq[A]\setminus Q(y)$ of cardinality $t$. Let $i\in S$. As $S\subseteq[A]\setminus Q(y)$ and $i\in S$ we have $i\notin Q(y)$ and therefore $B_i(y)=0$. Also, $y$ does not violate the product rule, hence,
\[(Y'''_i(y),Y'''_{S\setminus\{i\}}(y))\approx_{p_2}U_1\times Y'''_{S\setminus\{i\}}(y).\]
As this is true for any $i\in S$, by Lemma 18,
\[Y'''_S(y)\approx_{tp_2}U_t.\]
Thus, $Y'''(y)$ is a $(q,t,tp_2)$ non-oblivious bit-fixing source.
\end{proof}
In particular, by Lemma 20, for every good $y$, $Y'''(y)$ is $tA^tp_2$-close to a $(q,t)$ non-oblivious bit-fixing source. By the choices we have made above $q=\sqrt{\varepsilon_{\rm cond}}A\le A^{1-2\gamma}$. Equation (1) implies that
\[t=\frac{\delta_0}4\left(\frac1{\varepsilon_{\rm cond}}\right)^{\log(3/2)}\ge c_\gamma\log^4 A.\]
Using the resiliency of $f$ from Theorem 22 (and the fact that it is almost balanced), the output when $y$ is good is $p_3$-close to uniform for $p_3=tA^tp_2+\varepsilon_f+A^{-1/c_\gamma}$, where the first term is due to the distance from a $t$-wise distribution, the second is due to the resiliency and the third is due to the bias of $f$ (see, e.g., Lemma 2.11 in [9]). To that we also have to add the probability $p_1$ that $y$ is not good. To finish the proof we notice that:
\begin{itemize}
\item It holds that
\begin{align*}
\varepsilon_f&\le c_\gamma\frac q{A^{1-\gamma}}\le c_\gamma\frac{A^{1-2\gamma}}{A^{1-\gamma}}=c_\gamma A^{-\gamma}\\
&\le c_\gamma2^{-\gamma\log\frac1{\varepsilon_{\rm cond}}}\le\frac{\varepsilon_0}5,
\end{align*}
because $\varepsilon_{\rm cond}\le(\frac{\varepsilon_0}{5c_\gamma})^{1/\gamma}$.
\item Also,
\[A^{-1/c_\gamma}\le2^{-\frac1{c_\gamma}\log\frac1{\varepsilon_{\rm cond}}}=\varepsilon_{\rm cond}^{1/c_\gamma}\le\frac{\varepsilon_0}5,\]
because $\varepsilon_{\rm cond}\le(\frac{\varepsilon_0}5)^{c_\gamma}$.
\item Finally, $tp_2=t(\varepsilon_{\rm Raz}+2\varepsilon_{\rm CBA})$, $\varepsilon_{\rm Raz}=2^{-\Omega(k)}=d^{-\Omega(1)}$, $\varepsilon_{\rm CBA}=\frac1d$. Thus, $tp_2\le4td^{-\Omega(1)}$. $A$ and $t$ are constants, so for $d$ large enough, $tA^tp_2\le\frac{\varepsilon_0}5$.
\end{itemize}
Together, the error is at most $p_1+p_3\le\varepsilon_0$ completing the proof of the theorem.\hfill$\square$

\section*{Strong Seeded Dispersers and Friends (original Sections 5.1 and 5.3)}
\subsection*{5.1 Strong Seeded Dispersers}
\begin{definition}[strong disperser]
A function $\Disp:\{0,1\}^n\times\{0,1\}^d\to\{0,1\}^m$ is a strong $(k,\varepsilon)$ disperser, if for every $(n,k)$ source $X$,
\[|\Supp((Y,\Disp(X,Y)))|>(1-\varepsilon)DM.\]
We say $\Disp$ is (source) linear if for every $y\in\{0,1\}^d$ and every $x_1,x_2\in\mathbb F_2^n$, $\Disp(x_1+x_2,y)=\Disp(x_1,y)+\Disp(x_2,y)$.
\end{definition}
We are interested in the important special case where $m=1$. In this case, non-explicitly, a random function is (w.h.p.) a strong $(k,\varepsilon)$ disperser with $d=\log n+\log(\frac1\varepsilon)+O(1)$ provided that $k\ge\log\log(\frac1\varepsilon)+O(1)$ [37,33]. A matching lower bound, up to additive constant factors, was given by [37].
Using the translation between strong seeded dispersers and erasure list-decodable codes which we discuss in Section 5.3, Guruswami and Indyk's result [26] gives a probabilistic polynomial time algorithm that outputs with high probability a strong seeded disperser with seed-length $d=2\log(\frac1\varepsilon)+\log n+\log\log(\frac1\varepsilon)$ and optimal entropy-loss. The construction can be made deterministic, but with running time exponential in $1/\varepsilon$.
Note that as we discuss the one output bit case, the required entropy is essentially the entropy-loss. From Theorem 27 we can derive a better explicit construction of a strong disperser with small error.
\begin{theorem}
For every constant $0<\gamma<1$ there exists a constant $c\ge1$ such that for every integer $n$ and $\varepsilon\le n^{-\frac c{1-\gamma}}$ there exists an explicit strong $(k,\varepsilon)$ disperser $\Disp:\{0,1\}^n\times\{0,1\}^d\to\{0,1\}$ where $d=(1+\gamma)\log(\frac1\varepsilon)$ and $k=c\log d$.
\end{theorem}
\begin{proof}
Set $\varepsilon_0=\frac14$ and $\delta_0=\frac\gamma{1+\gamma}$. Let $c$ be the constant from Theorem 27 for $\delta_0$ and $\varepsilon_0$ and let $\twoExt:[N]\times[D]\to\{0,1\}$ be the $((n,k),(d,k_2=\delta_0d),\varepsilon_0)$ two-source extractor where $d=(1+\gamma)\log(\frac1\varepsilon)$ and $k=c\log d$. Notice that $d\ge c\log n$ (because $\varepsilon\le n^{-\frac c{1+\gamma}}$) as required. Let $\Disp(x,y)=\twoExt(x,y)$.
Let $X\subseteq[N]$ be a set of size $K$ and call a value $y\in[D]$ $b$-bad if $\Disp(X,y)=\{b\}$. It follows that the sets of 0-bad $y$-s and 1-bad $y$-s are each of size less than $K_2$. Therefore,
\begin{align*}
|\Supp((U_d,\Disp(X,U_d)))|&>2K_2+2(D-2K_2)=2D-2K_2\\
&=\left(1-\frac{K_2}D\right)2D=(1-\varepsilon)2D,
\end{align*}
because $\frac{K_2}D=2^{-(1-\delta_0)d}=2^{-\frac1{1+\gamma}d}=2^{-\log(\frac1\varepsilon)}=\varepsilon$.
\end{proof}
\subsection*{5.3 Erasure List-Decodable Codes}
An $(\bar n,n)$ (binary) code is a mapping $\mathcal C:\{0,1\}^n\to\{0,1\}^{\bar n}$. The code $\mathcal C$ is linear if $\mathcal C$ is linear, and is denoted by $[\bar n,n]$. We identify a code with the image of $\mathcal C$. For a linear $\mathcal C$ this image is a linear subspace of $\mathbb F_2^{\bar n}$ of dimension $n$. A generator matrix for an $[\bar n,n]$ code $\mathcal C$ is any matrix whose columns form a basis for $\mathcal C$. In the erasures noise model, an adversarially chosen subset of the codeword's symbols are erased and the positions where erasures have occurred are known.
\begin{definition}[erasure list-decodable code]
A code $\mathcal C\subseteq\{0,1\}^{\bar n}$ is $(s,L)$ erasure list-decodable if for every $r\in\{0,1\}^{\bar n-s}$ and every set $T\subseteq[\bar n]$ of size $\bar n-s$,
\[|\{c\in\mathcal C\mid c|_T=r\}|<L,\]
where $c|_T$ denotes the projection of $c$ to the coordinates in $T$.
\end{definition}
The following folklore lemma (see, e.g., [24, Lemma 1]) gives an alternative characterization of linear erasure list-decodable codes.
\begin{lemma}
An $[\bar n,n]_2$ linear code $\mathcal C$ is $((1-\varepsilon)\bar n,L)$ erasure list-decodable if and only if its $\bar n\times n$ generator matrix $G$ has the property that every $\varepsilon\bar n\times n$ sub-matrix of $G$ has rank greater than $n-\log L$.
\end{lemma}
Non-explicitly, we have:
\begin{theorem}[{[24]}]
For every $n$ and $\varepsilon>0$, there exists an $(\bar n,n)$ binary code that is $((1-\varepsilon)\bar n,L)$-erasure list-decodable of rate $\frac n{\bar n}=\Omega(\varepsilon)$ and $L=O(\log(1/\varepsilon))$.
\end{theorem}
Guruswami [25] observed that strong dispersers can be used to construct erasure list-decodable codes. Here we complement his argument, and note that strong dispersers are equivalent to erasure list-decodable codes. Given a function $\Disp:[N]\times[D]\to\{0,1\}$, we consider the $(D,n)$ code $\mathcal C_{\Disp}:\{0,1\}^n\to\{0,1\}^D$ defined by $\mathcal C_{\Disp}(x)_i=\Disp(x,i)$. Note that the code is linear if and only if $\Disp$ is linear.
\begin{lemma}[following {[25, Lemma 12]}]
The function $\Disp:[N]\times[D]\to\{0,1\}$ is a strong $(k,\varepsilon)$ disperser if and only if $\mathcal C_{\Disp}$ is $((1-2\varepsilon)D,K)$ erasure list-decodable.
\end{lemma}
\begin{proof}
For one direction, assume $\Disp$ is a strong $(k,\varepsilon)$ disperser. We wish to prove that $\mathcal C_{\Disp}$ is $((1-2\varepsilon)D,K)$ erasure list-decodable. Let $T=\{t_1,\ldots,t_{2\varepsilon D}\}\subseteq[D]$ be an arbitrary set of size $2\varepsilon D$ and $r\in\{0,1\}^{2\varepsilon D}$ an arbitrary string. Let $X_{T,r}\subseteq\{0,1\}^n$ denote the set of all the messages $x$ for which $\mathcal C_{\Disp}(x)|_T=r$. Then,
\[|\Supp((U_d,\Disp(X_{T,r},U_d)))|\le|T|\cdot1+(D-|T|)\cdot2\le(1-\varepsilon)2D,\]
where the first inequality follows by considering seeds in $T$ and seeds in $[D]\setminus T$. For a seed $t_i\in T$ we have that $\Disp(X_{T,r},t_i)$ is fixed, hence each such seed contributes 1 to the support size. For any other seed $y$, the support size of $\Disp(X_{T,r},y)$ is at most 2. As $\Disp$ is a strong $(k,\varepsilon)$ disperser, we conclude that $|X_{T,r}|\le K$ as desired.
For the other direction assume $\Disp$ is a not a strong $(k,\varepsilon)$ disperser. Then, there exists a set $X\subseteq\{0,1\}^n$ such that $|X|\ge K$ and $|\Supp((U_d,\Disp(X,U_d)))|\le(1-2\varepsilon)2D$. Note that for every $y\in[D]$ we have $|\Supp(\Disp(X,y))|\in\{1,2\}$. Therefore, following the above calculation, there exists a set $T\subseteq D$ of size at least $2\varepsilon D$ such that for each $y\in T$, $|\Supp(\Disp(X,y))|=1$. But this means that for every $x\in X$, $\mathcal C_{\Disp}(x)|_T$ is the same (punctured) codeword. It follows that $\mathcal C_{\Disp}$ is not $((1-2\varepsilon)D,K)$ erasure list-decodable.
\end{proof}
\begin{corollary}
If $\Disp:\{0,1\}^n\times\{0,1\}^d\to\{0,1\}$ is a strong $(k,\varepsilon)$ disperser with seed-length $d=a_1\log n+a_2\log(\frac1\varepsilon)+a_3$ (for some $a_1\ge1,a_2\ge1$ and $a_3$) then $\mathcal C_{\Disp}$ is a $((1-2\varepsilon)D,K)$ erasure list-decodable code of rate $2^{-a_3}\cdot n^{1-a_1}\cdot\varepsilon^{a_2}$.
\end{corollary}
When $\varepsilon$ is much smaller than $\frac1n$ the dominant factor is determined by $a_2$. As we mentioned earlier (and as Guruswami also notes in [25]) previous explicit constructions for binary codes had $a_2\ge2$ (usually inherited from extractor constructions). Our construction is the first to get arbitrary close to $a_2=1$ and small list-size. Combining Corollary 45 and Theorem 40, we obtain:
\begin{theorem}
For every constant $0<\gamma<1$ there exists a constant $c\ge1$ such that for every integer $n$ and $\varepsilon\le n^{-\frac c{1-\gamma}}$ there exists an explicit code $\mathcal C:\{0,1\}^n\to\{0,1\}^{(\frac1\varepsilon)^{1+\gamma}}$ that is
\[\left((1-2\varepsilon)\left(\frac1\varepsilon\right)^{1+\gamma},\left((1+\gamma)\log\frac1\varepsilon\right)^c\right)\]
erasure list-decodable of rate $n\varepsilon^{1+\gamma}$.
\end{theorem}

\begin{thebibliography}{45}
\bibitem{ref1} HL Abbott. Lower bounds for some Ramsey numbers. Discrete Mathematics, 2(4):289–293, 1972.
\bibitem{ref2} Noga Alon. The Shannon capacity of a union. Combinatorica, 18(3):301–310, 1998.
\bibitem{ref3} Noga Alon, Oded Goldreich, and Yishay Mansour. Almost k-wise independence versus k-wise independence. Information Processing Letters, 88(3):107–110, 2003.
\bibitem{ref4} Boaz Barak. A simple explicit construction of an n˜o(log n)-Ramsey graph. arXiv preprint, 2006. arXiv:math/0601651.
\bibitem{ref5} Boaz Barak, Guy Kindler, Ronen Shaltiel, Benny Sudakov, and Avi Wigderson. Simulating independence: New constructions of condensers, Ramsey graphs, dispersers, and extractors. Journal of the ACM (JACM), 57(4):20, 2010.
\bibitem{ref6} Boaz Barak, Anup Rao, Ronen Shaltiel, and Avi Wigderson. 2-source dispersers for no(1) entropy, and Ramsey graphs beating the frankl-wilson construction. Annals of Mathematics, 176(3):1483–1544, 2012.
\bibitem{ref7} Avraham Ben-Aroya, Dean Doron, and Amnon Ta-Shma. An efficient reduction from two- source to nonmalleable extractors: achieving near-logarithmic min-entropy. SIAM Journal on Computing, pages STOC17–31–STOC17–49, 2019.
\bibitem{ref8} Jean Bourgain. More on the sum-product phenomenon in prime fields and its applications. International Journal of Number Theory, 1(01):1–32, 2005.
\bibitem{ref9} Eshan Chattopadhyay and David Zuckerman. Explicit two-source extractors and resilient functions. Annals of Mathematics, 189(3):653–705, 2019.
\bibitem{ref10} Benny Chor and Oded Goldreich. Unbiased bits from sources of weak randomness and probabilistic communication complexity. SIAM Journal on Computing, 17(2):230–261, 1988.
\bibitem{ref11} Fan RK Chung. A note on constructive methods for Ramsey numbers. Journal of Graph Theory, 5(1):109–113, 1981.
\bibitem{ref12} Gil Cohen. Local correlation breakers and applications to three-source extractors and mergers. SIAM Journal on Computing, 45(4):1297–1338, 2016.
\bibitem{ref13} Gil Cohen. Non-malleable extractors-new tools and improved constructions. In LIPIcs-Leibniz International Proceedings in Informatics, volume 50. Schloss Dagstuhl-Leibniz-Zentrum fuer Informatik, 2016.
\bibitem{ref14} Gil Cohen. Two-source dispersers for polylogarithmic entropy and improved Ramsey graphs. In Proceedings of the 48th Annual ACM Symposium on Theory of Computing (STOC 2016), pages 278–284, 2016.
\bibitem{ref15} Gil Cohen. Two-source extractors for quasi-logarithmic min-entropy and improved privacy amplification protocols. In Electronic Colloquium on Computational Complexity (ECCC), volume 23, page 114, 2016.
\bibitem{ref16} Zeev Dvir, Swastik Kopparty, Shubhangi Saraf, and Madhu Sudan. Extensions to the method of multiplicities, with applications to kakeya sets and mergers. SIAM Journal on Computing, 42(6):2305–2328, 2013.
\bibitem{ref17} Zeev Dvir and Avi Wigderson. Kakeya sets, new mergers, and old extractors. SIAM Journal on Computing, 40(3):778–792, 2011.
\bibitem{ref18} Paul Erdös. Some remarks on the theory of graphs. Bulletin of the American Mathematical Society, 53(4):292–294, 1947.
\bibitem{ref19} Peter Frankl. A constructive lower bound for Ramsey numbers. Ars Combinatoria, 3(297- 302):28, 1977.
\bibitem{ref20} Peter Frankl and Richard M. Wilson. Intersection theorems with geometric consequences. Combinatorica, 1(4):357–368, 1981.
\bibitem{ref21} Parikshit Gopalan. Constructing Ramsey graphs from boolean function representations. In Proceedings of the 21st Annual IEEE Conference on Computational Complexity (CCC 2006), pages 14–pp, 2006.
\bibitem{ref22} Ronen Gradwohl, Guy Kindler, Omer Reingold, and Amnon Ta-Shma. On the error parameter of dispersers. In Approximation, Randomization and Combinatorial Optimization. Algorithms and Techniques, pages 294–305. Springer, 2005.
\bibitem{ref23} Venkatesan Guruswami. List decoding of error-correcting codes. PhD thesis, Massachusetts Institute of Technology, 2001.
\bibitem{ref24} Venkatesan Guruswami. List decoding from erasures: Bounds and code constructions. IEEE Transactions on Information Theory, 49(11):2826–2833, 2003.
\bibitem{ref25} Venkatesan Guruswami. Better extractors for better codes? In Proceedings of the 36th Annual ACM Symposium on Theory of Computing (STOC 2004), pages 436–444, 2004.
\bibitem{ref26} Venkatesan Guruswami and Piotr Indyk. Near-optimal linear-time codes for unique decoding and new list-decodable codes over smaller alphabets. In Proceedings of the 34th Annual ACM Symposium on Theory of Computing (STOC 2002), pages 812–821, 2002.
\bibitem{ref27} Xin Li. Three-source extractors for polylogarithmic min-entropy. In Proceedings of the 56th Annual IEEE Symposium on Foundations of Computer Science (FOCS 2015), pages 863–882. IEEE, 2015.
\bibitem{ref28} Xin Li. Improved non-malleable extractors, non-malleable codes and independent source extractors. In Proceedings of the 49th Annual ACM Symposium on Theory of Computing (STOC 2017), pages 1144–1156, 2017.
\bibitem{ref29} Xin Li. Non-malleable extractors and non-malleable codes: Partially optimal constructions. In Proceedings of the 34th Computational Complexity Conference (CCC 2019). Schloss Dagstuhl- Leibniz-Zentrum fuer Informatik, 2019.
\bibitem{ref30} Shachar Lovett. Recent advances on the log-rank conjecture in communication complexity. Bulletin of EATCS, 1(112), 2014.
\bibitem{ref31} Shachar Lovett. Communication is bounded by root of rank. Journal of the ACM (JACM), 63(1):1, 2016.
\bibitem{ref32} Raghu Meka. Explicit resilient functions matching Ajtai-Linial. In Proceedings of the Twenty- Eighth Annual ACM-SIAM Symposium on Discrete Algorithms (SODA 2017), pages 1132–1148, 2017.
\bibitem{ref33} Raghu Meka, Omer Reingold, and Yuan Zhou. Deterministic coupon collection and better strong dispersers. In Approximation, Randomization, and Combinatorial Optimization. Al- gorithms and Techniques (APPROX/RANDOM 2014). Schloss Dagstuhl-Leibniz-Zentrum fuer Informatik, 2014.
\bibitem{ref34} Zs Nagy. A constructive estimation of the Ramsey numbers. Mat. Lapok, 23:301–302, 1975.
\bibitem{ref35} Moni Naor. Constructing Ramsey graphs from small probability spaces. IBM Research Report RJ, 8810, 1992.
\bibitem{ref36} Noam Nisan and Avi Wigderson. On rank vs. communication complexity. Combinatorica, 15(4):557–565, 1995.
\bibitem{ref37} Jaikumar Radhakrishnan and Amnon Ta-Shma. Bounds for dispersers, extractors, and depth-two superconcentrators. SIAM Journal on Discrete Mathematics, 13(1):2–24, 2000.
\bibitem{ref38} FP Ramsey. On a problem of formal logic. Proceedings of the London Mathematical Society, 2(1):264–286, 1930.
\bibitem{ref39} Ran Raz. Extractors with weak random seeds. In Proceedings of the 37th Annual ACM Symposium on Theory of Computing (STOC 2005), pages 11–20, 2005.
\bibitem{ref40} Ronen Shaltiel. Recent developments in explicit constructions of extractors. Bulletin of the EATCS, 77(67-95):10, 2002.
\bibitem{ref41} Ronen Shaltiel. An introduction to randomness extractors. In International Colloquium on Automata, Languages, and Programming, pages 21–41. Springer, 2011.
\bibitem{ref42} Luca Trevisan. Extractors and pseudorandom generators. Journal of the ACM (JACM), 48(4):860–879, 2001.
\bibitem{ref43} Avi Wigderson. Randomness extractors – applications and constructions. In IARCS Annual Conference on Foundations of Software Technology and Theoretical Computer Science. Schloss Dagstuhl-Leibniz-Zentrum für Informatik, 2009.
\bibitem{ref44} David Zuckerman. Simulating BPP using a general weak random source. Algorithmica, 16(4):367–391, 1996.
\bibitem{ref45} David Zuckerman. Linear degree extractors and the inapproximability of max clique and chromatic number. Theory of Computing, 3:103–128, 2007. CCC 2020
\end{thebibliography}
\end{document}
